% Created 2018-05-10 Thu 15:04
% Intended LaTeX compiler: pdflatex
\documentclass[11pt]{article}
\usepackage[utf8]{inputenc}
\usepackage[T1]{fontenc}
\usepackage{graphicx}
\usepackage{grffile}
\usepackage{longtable}
\usepackage{wrapfig}
\usepackage{rotating}
\usepackage[normalem]{ulem}
\usepackage{amsmath}
\usepackage{textcomp}
\usepackage{amssymb}
\usepackage{capt-of}
\usepackage{hyperref}
\usepackage[margin=0.65in]{geometry}
\usepackage[backend=biber, style=ieee, url=false]{biblatex}
\usepackage{float}
\usepackage[super,negative]{nth}
\usepackage[capitalise]{cleveref}
\usepackage{pst-node,transparent,ragged2e}
\addbibresource{/home/nuk3/.spacemacs.d/org-files/bibliography.bib}
\DeclareFieldFormat[inproceedings]{citetitle}{\textit{#1}}
\DeclareFieldFormat[inproceedings]{title}{\textit{#1}}
\DeclareFieldFormat[misc]{citetitle}{#1}
\DeclareFieldFormat[misc]{title}{#1}
\renewcommand*{\bibpagespunct}{%
\ifentrytype{inproceedings}
{\addspace}
{\addcomma\space}}
\AtEveryCitekey{\ifuseauthor{}{\clearname{author}}}
\AtEveryBibitem{\ifuseauthor{}{\clearname{author}}}
\SpecialCoor
\author{Nyameko Lisa}
\date{13 April 2018}
\title{IPSP066 - Internet Aspects of Copyright and Trade Marks\\\medskip
\large Assignment 01 - 695242}
\hypersetup{
 pdfauthor={Nyameko Lisa},
 pdftitle={IPSP066 - Internet Aspects of Copyright and Trade Marks},
 pdfkeywords={},
 pdfsubject={},
 pdfcreator={Emacs 27.0.50 (Org mode 9.1.2)},
 pdflang={English}}
\begin{document}

\maketitle
\addvspace{110pt}
\centering{
\pnode(0.5\textwidth,-0.5\textheight){thisCenter}
\rput(thisCenter){%\transparent{0.25}
\includegraphics[width=2.7in]{/home/nuk3/course/llb/wipo-unisa/UNISACoatofArms.eps}}}

\justifying
\addvspace{110pt}
\section*{Declaration}
\label{sec:org26d54c1}
I know that plagiarism is to use someone else’s work and pass it off as my own.\\
I know that plagiarism is wrong.\\
I confirm that this assignment is my own work.\\
I have acknowledged in the bibliography accompanying the assignment all the sources that I have used.\\
I have not directly copied without acknowledgement anything from the Internet or from any other source.\\
I have indicated every quotation and citation in a footnote or bracket linked to that quotation.\\
I have not allowed anyone else to copy my work and to pass it off as their own work.\\
I understand that if any unacknowledged copying whatsoever appears in my assignment I will receive zero per cent for the assignment.\\
I am aware of the UNISA policy on plagiarism and understand that disciplinary proceedings can be instituted against me by UNISA if I contravene this policy.\\
I indicate my understanding and acceptance of this declaration by
entering my name here under:
\begin{itemize}
\item Name: \textbf{Nyameko Lisa} (Student Number: \textbf{7874-909-3})
\end{itemize}

\subsection*{NOTE}
\label{sec:orgd1eb897}
Please note that footnotes will be denoted as \footnote{This is a footnote.} and will
appear at the bottom of the page.\\
References will be denoted by \cite{rsa93_tm_act} and will appear at the end of the document.
\newpage

\section{Sibusiso wants to register the domain name `united'.}
\label{sec:orgf410807}

As per the definitions of [section 2(1)]\cite{rsa93_tm_act} and given that the
Republic is a Paris Convention Member State, [article 15(1)]\cite{wto17_trips},
[article 1(2)]\cite{wipo83_paris_conve_protect_ip}, it follows that
\emph{`UNITED.gTLD'} and \emph{`UNITED.ccTLD} will be understood to be \textbf{marks}, in that
they are names or signs capable of being represented graphically. Moreover
\emph{`UNITED.COM'} constitutes a \textbf{well-known registered service} or \textbf{trade mark} in
the United States of America in relation to distinguishing their trade or use
from other goods or services connected in the course of trade regarding
\emph{`airline and courier services.'}.

\subsection{Can he register such a domain name?}
\label{sec:org746194b}

In order to reserve such a domain name in a general Top-Level Domain or \textbf{gTLD},
the \textbf{registrant} Sibusiso must register it with an ICANN-accredited
registrar\footnote{Or alternatively through a registrar's resellers.}
\cite{icann13_rules}. The registrar will then check if the UNITED.gTLD domain name
is available, and then either create or lookup a WHOIS (domain name registration
record) with the registrant's information \cite{icann99_policy}. The registrant would be advised
against the pursuit of gTLD of the `united' domain name, as he is not likely to
be successful. In particular, while the `united' trade mark may be territorial by
its very nature, the specific `united.com' domain name cannot be owned by
multiple companies, irrespective of the uniqueness of their products and / or
markets.\\

Alternatively, the registrant would be strongly advised to instead pursue
registration of the country code Top-Level Domain or \textbf{ccTLD} domain
`united.co.za', through one of the ccTLD operators
[section 59]\cite{rsa02_elect_comm_trans_act}, particularly if the students he intends to sell
sound recordings to are based within the Republic. This will be assigned on a
first-come first-serve basis.

\subsection{Could he be liable for trade mark infringement?}
\label{sec:org9166777}
Yes indeed, as per the provisions of [sections 10(6) and 35]\cite{rsa93_tm_act}
\emph{(for litigation within the Republic)}, [Article 16 subsections (1), (2) and
(3)]\cite{wto17_trips}, [Article
6\textsuperscript{\textit{bis}}(1)]\cite{wipo83_paris_conve_protect_ip},
[Article 16]\cite{wipo70_tlt} and [subsection 1125(c)]\cite{usa46_title15_ch22}
\emph{(for actions within the United States of America)}, the proprietors of the
second level domain `united', as is the case with the airline company, may argue
that it is indeed a \textbf{well-known} mark and that registration of Sibusiso's domain
name in either a gTLD or ccTLD, constitutes trade mark infringement, through the
\textbf{dilution} of a well-known mark.\\

Dilution protection is not subject to the same limitations of `traditional'
trade mark infringement,\footnote{Restricted to goods and services which are the same or similar to those for which the mark is registered.\label{orgc5d7a2b}} in that it extends to use of the offending mark to
\emph{any} goods and services. The proprietors of the UNITED.com Airlines trade
mark and domain name, need demonstrate that:
\begin{itemize}
\item Their mark is well known within the Republic, i.e. it would be sufficient for
them to show that their trade mark has acquired a reputation amongst a
substantial number of members of the public,
\citetitle{corbett97_mcd_v_joburgers} \cite{corbett97_mcd_v_joburgers},
\item Sibusiso has used the mark in the course of trade, where the courts of
the Republic view this in the `traditional' trade mark infringement\textsuperscript{\ref{orgc5d7a2b}}
sense, \citetitle{harms01_cowbell_v_ics} \cite{harms01_cowbell_v_ics}, lastly
\item Use of the infringing mark is likely to take unfair advantage of, or be
detrimental to the distinctive character of the UNITED Airlines trade mark,
amounting to unfair competition [articles 2(1) and 3(2)(a)(i)]\cite{wipo96_model_provi_unfair_comp}.
\end{itemize}

\subsection{Is there any other procedure that could be instituted to resolve the dispute?}
\label{sec:orgb1c6127}

ICANN provides a Uniform Domain Name Dispute Resolution Policy which outlines
provisions [section 4(a)(i)]\cite{icann99_policy}, within which mandatory
administrative proceedings may be instituted, on the basis that as a \textbf{third
party complainant} UNITED Airlines asserts that Sibusiso's domain is identical
or confusingly similar to their trade or service mark over which they have
rights to.

\section{Sibusiso wants to place samples from the sound recordings for sale on his web site.}
\label{sec:org6b4d04c}
As a Convention Member the provisions of [Article 9(1)]\cite{wto17_trips},
[section 5]\cite{rsa78_copyrightact}, [Article 1(2)]\cite{eurlex00_elec_commerce},
[Article 5(2)]\cite{wipo86_berne} and [Article 1(4)]\cite{wipo96_copyright_treaty}
which dictate that the statues of the Republic shall govern international
copyright protection as it pertains to the Internet. Furthermore as per the
provisions of [sections 1 and 2(1)(b,e)]\cite{rsa78_copyrightact} and [Article
2(1)]\cite{wipo86_berne} the sound recordings shall be considered as \textbf{protected
artistic works} or \textbf{sound recordings}.

\subsection{Which holders of intellectual property rights may be affected by this?}
\label{sec:orgcf47608}

As per the provisions of [Article 14]\cite{wto17_trips}, [sections 1, 3,
4(1)(a)]\cite{rsa78_copyrightact}, [Articles 2 and 3]\cite{wipo96_wppt}, [Article 3
and 9(1)]\cite{wipo86_berne}, the \textbf{authors}, and they are understood to be the
performers, producers, licensees and broadcasting organisations of the sound
recordings.

\subsection{Which of their rights may be affected?}
\label{sec:org1bfd797}
\begin{itemize}
\item Exclusive right of reproduction [section 9(a)]\cite{rsa78_copyrightact},
[Article 9(1) and (2)]\cite{wipo86_berne}, [Article 7 and 11]\cite{wipo96_wppt}
and [Article 2(a-e)]\cite{eurlex01_harmonisation},

\item Exclusive right of offering by way of trade, directly or indirectly, a
reproduction of the sound recording [section 9(b)]\cite{rsa78_copyrightact},

\item Exclusive right of publication, distribution or broadcast to the public, or
causing the sound recording to be transmitted in a diffusion service [section
9(c-d)]\cite{rsa78_copyrightact}, [Article
11\textsuperscript{bis}(1)(i)-(ii)]\cite{wipo86_berne}, [Article
6(1)]\cite{wipo96_copyright_treaty}, [Article 8(1) and 12(1)]\cite{wipo96_wppt}
and [Article 4]\cite{eurlex01_harmonisation},

\item Exclusive right of communication to the public of a performance [section
9(e)]\cite{rsa78_copyrightact}, [Article 11(1)(i)-(ii)]\cite{wipo86_berne},
[Article 8]\cite{wipo96_copyright_treaty}, [Article 13-15]\cite{wipo96_wppt} and
[Article 3]\cite{eurlex01_harmonisation}, and

\item Exclusive right\footnote{Sampling constitutes an adaptation.} of adaptation,
arrangement and other alteration [section 6(f-g)]\cite{rsa78_copyrightact},
[Article 12]\cite{wipo86_berne} and [Article 16]\cite{wipo96_wppt}.
\end{itemize}

\subsection{Does he need their permission to do so? What form will such permission take?}
\label{sec:orgc0aa21d}

Yes indeed, by way of a license, [section 22]\cite{rsa78_copyrightact} and
[Article 13(1)]\cite{wipo86_berne}.

\subsection{If he does not obtain their permission, will HostNet also be legally liable?}
\label{sec:org383c837}

With respect to Sibusiso's intellectual property rights infringement, the
intermediary service provider HostNet shall \uline{\textbf{not}} be legally liable, for
activities conducted during the course of its trade:
\begin{itemize}
\item Information transmission or acting as a \textbf{`mere conduit'} [Article 12]\cite{eurlex00_elec_commerce},
\item Temporary information storage \textbf{`caching'} [Article 13]\cite{eurlex00_elec_commerce}, and
\item Hosting [Article 14]\cite{eurlex00_elec_commerce}.
\end{itemize}

Moreover there exist provisions against statutory requirements obliging HostNet
to monitor Sibusiso's activities on their respective servers [Article 15]\cite{eurlex00_elec_commerce}.

\printbibliography
\end{document}